\section{The density-free foundation}
\label{sec:geometry}

We record the precise geometric input from
\cite{Foundation2026}. Every statement in this section is
density-free and uses only the fixed compact manifold and weak
MTW. The proofs, including the treatment of conjugate cut
endpoints, are in that paper.

A potential is a function
$u(x)=\max_y\{-c(x,y)-w(y)\}$ with a continuous datum $w$;
replacing $w$ by the cost transform of $u$ gives a dual pair.
Put $D_g=\operatorname{diam}M$.
The total minimizing-vector set $\{(x,p):p\in G_x\}$ is compact.
For $p\in G_x$ and $y_s=\exp_x(sp)$, $0<s<1$, divided
action gives
\begin{equation}\label{geo:split-support}
 c(z,\exp_xp)\leq c_s(z,y_s)+\frac{1-s}{2}|p|^2,
 \qquad c_s(x,y_s)=\frac{s}{2}|p|^2.
\end{equation}
The right side is a smooth upper support near $x$.

\begin{lemma}[Uniform supports and active logs]\label{geo:prelim}
There is a geometric constant $C_g$ such that the cost is
uniformly semiconcave in either variable, including at cut
points. Every potential is $D_g$-Lipschitz and uniformly
semiconvex. For a continuous representing datum $w$, define
\[
 \mathcal V_u(x)=\{p\in G_x:
              u(x)+c(x,\exp_xp)+w(\exp_xp)=0\}.
\]
Then $\mathcal V_u(x)$ is compact and
$\partial u(x)=\operatorname{conv}\mathcal V_u(x)$.
Every $p\in\partial u(x)$ has a positive representation using at most
$n+1$ affinely independent active logs.
The set $\mathcal V_u(x)$ includes every minimizing log to
each endpoint active for the specified datum $w$.
\end{lemma}

This is \cite[Lemma~4.1]{Foundation2026}. In particular, the smooth upper
supports obtained by splitting at time $1/2$ have a common
neighborhood size and Hessian bound. The representing datum
need not itself be a cost potential.

Write $H_s(p)=A_x(sp)/s$. At a specified nonconjugate
minimizing branch, $A_x$ will also denote the smooth branch
Hessian.

\begin{lemma}[Radial divided action]\label{geo:radial-index}
Let $p\in G_x$, and let $S_p(r)$ be the Jacobi tensor along
$\exp_x(rp)$ with $S_p(0)=0$, $D_rS_p(0)=\Id$, in parallel
orthonormal frames. For $0<t<s<1$,
\[
 \frac{d}{dr}H_r(p)=-S_p(r)^{-1}S_p(r)^{-T},
 \qquad H_t(p)-H_s(p)\geq\kappa_g(s-t)\Id,
\]
where $\kappa_g>0$ depends only on the manifold. On every
smooth minimizing branch, $A_x(w)w=w$.
\end{lemma}

\begin{lemma}[Transverse slack]\label{geo:slack}
Let $p_i\in G_x$, $m_i>0$, $\sum_i m_i=1$, and set
$p=\sum_i m_ip_i$, $E=\Span\{p_i-p\}$. If
$t\operatorname{conv}\{p_i\}\subset I(x)$ and $0<t<s<1$,
then
\[
 \left[\frac{A_x(tp)}t-\sum_i m_i\frac{A_x(sp_i)}s\right]
       \bigg|_{E^\perp}\geq\kappa_g(s-t)\Id.
\]
The conclusion also holds at a limiting time $t$ if the
radially contracted hulls lie in $I(x)$ and the displayed
Hessians have nonconjugate branch limits. No minimizing
assumption about $sp$ is needed in this limiting assertion.
\end{lemma}

These are \cite[Lemmas~4.2--4.3]{Foundation2026}. The latter applies to
every positive finite representation, not merely a selected
Carath\'eodory representation. We reconstruct the Jacobi
identities and the additional uniform spectral bounds in
Section~\ref{sec:spectral}.

\begin{lemma}[Short time for a semiconvex datum]\label{geo:short-time}
Let $f$ be Lipschitz and semiconvex. For sufficiently small
$t>0$, depending only on these two bounds and the manifold,
\[
 F_t:\Gamma_f\to M,\qquad
 \Gamma_f=\{(x,p):p\in\partial f(x)\},\qquad
 F_t(x,p)=\exp_x(tp)
\]
is a homeomorphism with locally Lipschitz inverse into the
tangent bundle. The minimizing pole $x_t(z)$ of $Q_tf$ is
unique, $Q_tf\in C^{1,1}$, and
$\grad Q_tf(z)=-t^{-1}\log_zx_t(z)$.
The graph $\Gamma_f$ is an $n$-dimensional Lipschitz manifold.
\end{lemma}

This is \cite[Lemma~4.4]{Foundation2026}. The local mechanism is the curved
version of proximal regularization and quadratic
inf-convolution; compare \cite[Sections~3--5 and~7]{Moreau1965}.
It is needed below for increasing scalar modifications of
potentials, which need not themselves be cost potentials.

\begin{theorem}[Global supporting property]\label{geo:support}
For every potential $u$, every $x\in M$, and every
$p\in\partial u(x)$, one has $p\in G_x$ and
\[
 u(x')\geq u(x)+c(x,\exp_xp)-c(x',\exp_xp)
 \qquad(x'\in M).
\]
For each $0<t<1$, $F_t:\Gamma_u\to M$ is a
homeomorphism, $tp\in I(x)$, and its poles are the unique
global minimizers of $Q_tu$.
\end{theorem}

\begin{corollary}[Convex injectivity domains]\label{geo:convexity}
Both $G_x$ and $I(x)$ are convex for every $x\in M$.
\end{corollary}

\begin{corollary}[Convex lifted gap sections]\label{geo:sections}
For a dual pair $(u,v)$, each set
\[
 \mathcal S_x(r)=\left\{p\in G_x:
      u(x)+\frac{|p|^2}{2}+v(\exp_xp)\leq r\right\}
\]
is compact and convex, including ordinary and conjugate cut
endpoints. If it is nonempty and $r\geq0$, then
\[
 \operatorname{diam}\mathcal S_x(r)^2
 \leq8\left(r+\osc_{p\in\mathcal S_x(r)}v(\exp_xp)\right).
\]
\end{corollary}

These are \cite[Theorem~4.8 and Corollaries~4.9--4.10]{Foundation2026}.
Convexity of
injectivity domains alone is not the whole input: the
ordinary-subgradient supporting assertion is also used for
modified envelopes and local lower tests.

\begin{proposition}[Intermediate regularity and pole control]
\label{geo:intermediate}
Let $(u,v)$ be a dual pair and $0<t<1$. Then
$Q_tu=-Q_{1-t}v\in C^{1,1}(M)$, with almost-everywhere bounds
\[
 -\frac{C_g}{1-t}g\leq\Hess Q_tu\leq\frac{C_g}{t}g.
\]
If $(x_t(z),p_t(z))=F_t^{-1}(z)$ and $\Phi_s$ is geodesic
flow on $TM$, then
\[
 (x_t(z),p_t(z))=\Phi_{-t}(z,\grad Q_tu(z)).
\]
In particular, the pole and momentum maps are Lipschitz
with constants depending only on $M,g,t$.
There are also $r_t,b_t>0$, depending only on $M,g,t$,
such that every $z=\exp_x(tp)$, $p\in\partial u(x)$,
satisfies
\[
 u(x')+c_t(x',z)\geq
 u(x)+c_t(x,z)+b_t d(x,x')^2
 \qquad\text{if }d(x,x')<r_t.
\]
All these constants are uniform when $t$ ranges in a
compact subinterval of $(0,1)$.
\end{proposition}

This is \cite[Proposition~4.11]{Foundation2026}. Its proof uses the two-sided
supports from the split contact geodesic and smooth
backward geodesic flow, not inversion of the full
endpoint exponential map.

\begin{remark}\label{geo:time-one}
No time-one injectivity is claimed for an arbitrary
potential. For example, $u(x)=-c(x,y_0)$ can have many
poles with endpoint $y_0$. All inverse projections and
$C^{1,1}$ regularizations in this paper are used at
times strictly between zero and one.
\end{remark}
